\chapter{Conclusions}
\label{cap:conclusions}

% TARGET LENGTH: ~4 pages.

\section{Summary of Contributions}
\label{sec:concl-summary}

% Recap the four contributions stated in Chapter 1; tie each one to
% concrete deliverables (architecture, codebase, tests, documentation).

\section{Goals Achieved}
\label{sec:concl-goals}

% Walk through the seven objectives from Section 1.3 with a one-line
% verdict per objective and a pointer to the chapter where it was
% delivered.

\section{Lessons Learned}
\label{sec:concl-lessons}

\subsection{Product Owner Lessons in a Two-Person Team}
% Discuss the tension of being PO and main implementer. How backlog
% ordering was kept client-driven. How sprint reviews became the main
% mechanism for the teammate to push back on PO decisions. The risk of
% scope creep when the same person holds both roles.

\subsection{Architectural Lessons}
% The onion rule paid off the moment the AI Interviewer had to be
% extracted into a separate process. The cost of "no ORM" (manual
% camelize/snakeify) was real but bounded; type safety filled the gap.

\subsection{AI Lessons}
% Structured-output prompting beats free-form for any field that has to
% land in a database. Schema versioning of LLM outputs is a must.
% Fallback providers are a reliability tool, not a luxury.

\section{Limitations}
\label{sec:concl-limitations}

% Sample-size limits of the validation survey; coverage gaps in some
% use-case modules; reliance on browser-native STT/TTS; manual
% right-to-erasure flow.

\section{Future Work}
\label{sec:concl-future}

% Backlog from docs/audit-6-5-2026.md items #2--#14; right-to-erasure
% automation; expansion of the analytics widget catalogue; offline
% transcription via Whisper for browsers without Web Speech API;
% multi-tenant separation (organisation-level data isolation).

\section{Closing Remarks}
\label{sec:concl-closing}

% Brief, sober closing paragraph.
